\documentclass[UTF8,fontset=fandol,a4paper,10pt]{ctexart}

\usepackage[margin=18mm]{geometry}
\usepackage{amsmath,amssymb,mathtools,booktabs,array}
\usepackage[hidelinks]{hyperref}

\setlength{\parindent}{0pt}
\setlength{\parskip}{3pt}
\setlength{\abovedisplayskip}{6pt}
\setlength{\belowdisplayskip}{6pt}
\ctexset{section={format=\large\bfseries,beforeskip=9pt,afterskip=4pt}}

\newcommand{\BF}{\mathrm{BF16}}
\newcommand{\FP}{\mathrm{FP32}}
\newcommand{\RN}{\operatorname{R}_{\BF}}
\newcommand{\Qnt}{\mathcal{Q}_{16}}
\newcommand{\N}{\mathrm{N}}
\newcommand{\R}{\mathrm{R}}
\newcommand{\diag}{\operatorname{diag}}
\newcommand{\CTA}{\mathrm{CTA}}
\newcommand{\lse}{\operatorname{LSE}}

\begin{document}

\begin{center}
  {\Large\bfseries SM120 NVFP4 Sparse Attention：数学数据流}\\[3pt]
  {\small 当前 grouped kernel；$H=64$，$D=448+64$，每个 query 单独处理}
\end{center}

以下显式保留 NVFP4 量化与 BF16 舍入；MMA 累加、softmax 状态及合并权重使用 FP32。
矩阵式描述其计算与数据依赖，不展开每条 FP32 指令的舍入。

\section*{1\quad 输入、稀疏选取与 CTA 划分}

每个请求只有一个 query token。输入与输出形状为
\[
  Q,O\in\BF^{B\times1\times64\times512},\qquad
  B\in\{64,256,1024\},\qquad \eta=\frac1{\sqrt{512}}.
\]
固定请求 $b$，省略其下标；$\N$ 表示 non-RoPE，$\R$ 表示 RoPE：
\[
  Q=[Q^{\N}\mid Q^{\R}],\qquad
  Q^{\N}\in\BF^{64\times448},\quad Q^{\R}\in\BF^{64\times64}.
\]
两个索引序列分别从 SWA 缓存选取 128 项、从压缩缓存选取 512 项。
历史长度 $32768$、压缩率 $4$ 时，后者从 $8192$ 个缓存条目中选取：
\[
  I^{\mathrm{swa}}\in\mathbb Z^{128},\quad
  I^{\mathrm{cmp}}\in\mathbb Z^{512},\qquad
  C=\left\lceil\frac{128}{64}\right\rceil+
    \left\lceil\frac{512}{64}\right\rceil=10.
\]
对 chunk $c$、其中候选位置 $j\in\{0,\ldots,63\}$，实际 gather 为
\[
  K_{c,j}=
  \begin{cases}
    \mathrm{Cache}_{\mathrm{swa}}[I^{\mathrm{swa}}_{64c+j}],&c=0,1,\\
    \mathrm{Cache}_{\mathrm{cmp}}[I^{\mathrm{cmp}}_{64(c-2)+j}],&c=2,\ldots,9.
  \end{cases}
\]
上式表示解码缓存条目后的逻辑值。$K_c=[K_c^{\N}\mid K_c^{\R}]$ 的形状为
$64\times512$；同一 latent KV 条目同时提供 K 和 V，64 个 head 共用它。
页大小为 64，索引 $i$ 对应 $(\lfloor i/64\rfloor,\,i\bmod64)$。
超出有效长度或索引为负的候选置零并屏蔽；参与集合由输入索引指定。
每个 $\CTA_{b,s}$ 负责全部 64 个 head 及一段连续 chunk：
\[
  \mathcal C_s=\{sp,\ldots,\min((s+1)p,10)-1\},
  \qquad p=\text{chunks per block}.
\]
当前测得的两种划分为
\[
  \begin{array}{c|c|c}
    B&\text{每个请求的主计算 CTA}&\text{各 CTA 处理的 chunk}\\\hline
    64&2\quad(\mathrm{cpb}=9)&\{0,\ldots,8\},\ \{9\}\\
    256,1024&1\quad(\mathrm{cpb}=10)&\{0,\ldots,9\}
  \end{array}
\]
这些 CTA 可由不同 SM 调度执行，互不共享 softmax 状态；两片结果随后单独合并。

\section*{2\quad 量化与 QK}

对 16 个数 $x=(x_0,\ldots,x_{15})$，定义 kernel 使用的量化--重建算子
\[
  a=\max_i|x_i|,\qquad s=R_{\mathrm{E4M3}}(a/6),\qquad
  r=\begin{cases}1/s,&s\ne0,\\0,&s=0,\end{cases}
  \qquad \Qnt(x)_i=s\,R_{\mathrm{E2M1}}(x_i r).
\]
$R$ 表示最近偶数舍入，FP4/FP8 越界时作有限值饱和；E2M1 的幅值为
$\{0,\tfrac12,1,\tfrac32,2,3,4,6\}$。
$\Qnt$ 表示 payload 与 scale 所代表的数值，实际仍以 FP4 和 E4M3 存储。
Q 的 non-RoPE 部分按每行连续 16 维量化；K 直接使用缓存已有的 FP4 与 scale：
\begin{align*}
  \widehat Q^{\N}_{h,16g:16g+16}
    &=\Qnt(Q^{\N}_{h,16g:16g+16}),\qquad g=0,\ldots,27,\\
  K^{\N}_{c,j,d}
    &=s^K_{c,j,\lfloor d/16\rfloor}\,e^K_{c,j,d},\qquad d=0,\ldots,447.
\end{align*}
切片均为左闭右开。缓存的 RoPE 部分 $K_c^{\R}$ 及 $Q^{\R}$ 保持 BF16。
先累加 448 维 NVFP4 点积，再累加 64 维 BF16 点积：
\[
  Z_c=\eta\log_2e\,
  \Bigl(
    \underbrace{\widehat Q^{\N}(K_c^{\N})^\top}_{(64\times448)(448\times64)}
    +\underbrace{Q^{\R}(K_c^{\R})^\top}_{(64\times64)(64\times64)}
  \Bigr)\in\FP^{64\times64}.
\]
无效位置设 $Z_{c,hj}=-10^{30}$。这里不重新执行 RoPE 变换。

\section*{3\quad 单个 CTA 内：逐 chunk 的 online softmax 与 PV}

以下固定一个 CTA；它按自身 chunk 顺序循环。初始化
\[
  m_h=-10^{30},\qquad \ell_h=0,\qquad
  A^{\N}=0_{64\times448},\qquad A^{\R}=0_{64\times64}.
\]
$m,\ell$ 是累计最大值与指数和；$A$ 是尚未归一化的输出累加量。

\textbf{(a) 当前 chunk 的局部归约。}

计算 warp $w\in\{0,\ldots,7\}$ 负责候选集合
$J_w=\{8w,\ldots,8w+7\}$，并处理全部 64 个 head（每组 16 个）。
设 $v_{c,j}$ 表示候选有效，则
\begin{align*}
  \mu_{h,w}&=\max_{j\in J_w} Z_{c,hj},\\
  p^{\mathrm{loc}}_{hj}&=
    \begin{cases}2^{Z_{c,hj}-\mu_{h,w}},&v_{c,j}=1,\\0,&v_{c,j}=0,\end{cases}
    \qquad j\in J_w,\\
  \ell_{h,w}&=\sum_{j\in J_w}p^{\mathrm{loc}}_{hj},\qquad
  \mu_h=\max_w\mu_{h,w},\qquad
  L_h=\sum_{w=0}^{7}\ell_{h,w}\,2^{\mu_{h,w}-\mu_h}.
\end{align*}
每个 head 的 $\mu_h,L_h$ 覆盖当前 chunk 的全部 64 个候选。

\textbf{(b) 合并到已有 softmax 状态。}

以上一轮的 $m_h,\ell_h$ 为输入，计算
\begin{align*}
  m'_h&=\max(m_h,\mu_h),\qquad
  \alpha_h=\begin{cases}2^{m_h-m'_h},&m_h>-10^{29},\\0,&\text{否则},\end{cases}\\
  \ell'_h&=\alpha_h\ell_h+2^{\mu_h-m'_h}L_h,\\
  W_{hj}&=\RN\!\left(p^{\mathrm{loc}}_{hj}\,2^{\mu_{h,w}-m'_h}\right),
  \qquad j\in J_w,\quad W\in\BF^{64\times64}.
\end{align*}
$W$ 是当前最大值基准下的指数权重，\textbf{尚未除以 $\ell'$}。
分母 $\ell'$ 来自 FP32 指数和，不从量化后的权重重新求和。

\textbf{(c) 准备 PV 的两个 NVFP4 操作数。}

权重沿候选轴每 16 项量化；V 先吸收缓存 scale，再沿候选轴重新量化：
\begin{align*}
  \widehat P_{h,16t:16t+16}
    &=\Qnt(W_{h,16t:16t+16}),&&t=0,\ldots,3,\\
  \widehat V^{\N}_{16t:16t+16,d}
    &=\Qnt(K^{\N}_{c,16t:16t+16,d}),&&d=0,\ldots,447.
\end{align*}
因此 V 的新 scale 跨越 16 个候选，和缓存中跨越 16 个特征维度的 scale 不同。
其转置布局为 $(\widehat V^{\N})^\top\in\mathbb R^{448\times64}$。

\begin{center}
\begin{tabular}{lccc}
\toprule
操作数 & 逻辑形状 & 每个 scale 覆盖 & scale 形状\\
\midrule
$\widehat Q^{\N}$ & $64\times448$ & 同一 head 的 16 维 & $64\times28$\\
$K_c^{\N}$ & $64\times448$ & 同一候选的 16 维 & $64\times28$\\
$(\widehat V^{\N})^\top$ & $448\times64$ & 同一维度的 16 个候选 & $448\times4$\\
$\widehat P$ & $64\times64$ & 同一 head 的 16 个候选 & $64\times4$\\
\bottomrule
\end{tabular}
\end{center}

\textbf{(d) 重缩放旧输出，再累加当前 chunk 的 PV。}

\begin{align*}
  A^{\N\prime}
    &=\diag(\alpha)A^{\N}
      +\underbrace{\widehat P\widehat V^{\N}}_{(64\times64)(64\times448),\ \mathrm{NVFP4}},\\
  A^{\R\prime}
    &=\diag(\alpha)A^{\R}
      +\underbrace{W K_c^{\R}}_{(64\times64)(64\times64),\ \BF}.
\end{align*}
更新 $(m,\ell,A^{\N},A^{\R})\leftarrow(m',\ell',A^{\N\prime},A^{\R\prime})$，
再处理下一个 chunk。PV 阶段各计算 warp 分担输出维度；RoPE 部分每个 warp 负责 8 维。
同一 CTA 内，IO warp 的 $K_c\rightarrow(\widehat V_c^{\N})^\top$
与计算 warp 的 $QK_c\rightarrow W_c\rightarrow\widehat P_c$ 并行推进，
两路就绪后执行 NVFP4 PV。原始 KV 使用双缓冲，转换后的 V 使用单缓冲；
V 缓冲在 NVFP4 PV 后释放，原始 KV 缓冲在 BF16 PV 后释放。

\section*{4\quad CTA 输出、sink 与 split 合并}

循环结束后，令 $A_s=[A_s^{\N}\mid A_s^{\R}]\in\FP^{64\times512}$。
对有有效候选的 head，当前 split 的 base-2 LSE 为
\[
  \lambda_{s,h}=m_{s,h}+\log_2\ell_{s,h}.
\]
sink 输入为每个 head 的实数 $u_h$，设 $\beta_h=u_h\log_2e$。
它只向 softmax 分母增加 $e^{u_h}$，对应的输出向量为零。

\textbf{单 CTA 直接输出：当前 $B=256,1024$ 的路径。}

不先将 $A/\ell$ 舍入为 BF16，而是先算 sink 缩放再一次写回：
\begin{align*}
  M_h&=\max(\lambda_h,\beta_h),\qquad
  d_h=2^{\lambda_h-M_h}+2^{\beta_h-M_h},\\
  \rho_h&=\frac{2^{\lambda_h-M_h}}{d_h},\qquad
  \boxed{\ O_{h,:}=\RN\!\left(\frac{\rho_h}{\ell_h}A_{h,:}\right)\ },\\
  \lse_h&=\kappa\bigl(M_h+\log_2 d_h\bigr).
\end{align*}
没有 sink 时取 $\rho_h=1$，最终 LSE 为 $\kappa\lambda_h$。

\textbf{两个 CTA 分片输出：当前 $B=64,\mathrm{cpb}=9$ 的路径。}

两个主计算 CTA 分别处理 $9\times64=576$ 与 $1\times64=64$ 个候选。
每片先保存已归一化的 BF16 中间结果，此时不加入 sink：
\[
  \widetilde O_{s,h,:}=\RN\!\left(\frac{A_{s,h,:}}{\ell_{s,h}}\right),\quad s=0,1;
  \qquad
  \widetilde O\in\BF^{B\times64\times2\times512},\quad
  \lambda\in\FP^{B\times64\times2}.
\]
合并 kernel 为每个请求分配 4 个 CTA，各负责 16 个 head；对每个 head 计算
\begin{align*}
  M_h&=\max(\lambda_{0,h},\lambda_{1,h},\beta_h),\\
  e_{s,h}&=2^{\lambda_{s,h}-M_h},\qquad e_{\mathrm{sink},h}=2^{\beta_h-M_h},\\
  d_h&=e_{0,h}+e_{1,h}+e_{\mathrm{sink},h},\qquad
  \omega_{s,h}=\frac{e_{s,h}}{d_h},\\
  \lse_h&=\kappa\bigl(M_h+\log_2d_h\bigr).
\end{align*}
\[
  \boxed{\ O_{h,:}=\RN\!\left(\omega_{0,h}\widetilde O_{0,h,:}
                  +\omega_{1,h}\widetilde O_{1,h,:}\right)\ }.
\]
没有 sink 时去掉其最大值项，并取 $e_{\mathrm{sink},h}=0$。
中间 LSE 始终以 2 为底；$\kappa$ 为最终 LSE 的缩放参数，
$\kappa=1$ 保留 base-2，$\kappa=\ln2$ 转成自然对数。
空 split 保存零输出及 $\lambda=-10^{30}$，合并时对应 $e_s=0$。
若整个请求没有有效候选，则 $O=0$；有 sink 时 $\lse=\kappa\beta$，否则为 $-\infty$。


\section*{5\quad Online Softmax 更新与分块 PV 的数学等价性}

以下固定一个 head $i$，省略所有 head 下标。令当前 chunk 中候选 $j$ 的
score 简记为
\[
  z_j:=Z_{c,ij}.
\]
由于第 2 节已经将 $\log_2e$ 乘入 $Z$，本节统一使用 $2^x$ 计算指数，
并暂时忽略量化和舍入误差。

\textbf{(a) 从 warp 局部状态得到当前 chunk 的状态。}

当前 chunk 的 64 个候选被划分给 8 个 warp：
\[
  J_w=\{8w,\ldots,8w+7\},\qquad w=0,\ldots,7.
\]
warp $w$ 计算
\begin{align*}
  \mu_w
    &=\max_{j\in J_w}z_j,\\
  p^{\mathrm{loc}}_j
    &=2^{z_j-\mu_w},\qquad j\in J_w,\\
  \ell_w
    &=\sum_{j\in J_w}p^{\mathrm{loc}}_j.
\end{align*}
无效候选的 $p^{\mathrm{loc}}_j$ 取零。

各 warp 使用不同的局部最大值 $\mu_w$，因此不能直接将
$\ell_w$ 相加。先计算整个 chunk 的最大值
\[
  \mu=\max_w\mu_w,
\]
再将每个 warp 的局部指数和转换到统一基准 $\mu$：
\[
  L
  =
  \sum_{w=0}^{7}\ell_w2^{\mu_w-\mu}.
\]
对每个 warp，有
\begin{align*}
  \ell_w2^{\mu_w-\mu}
  &=
  \left(
    \sum_{j\in J_w}2^{z_j-\mu_w}
  \right)
  2^{\mu_w-\mu}\\
  &=
  \sum_{j\in J_w}2^{z_j-\mu}.
\end{align*}
因此
\[
  \boxed{
    L=\sum_{j=0}^{63}2^{z_j-\mu}
  }.
\]
其中无效候选对应项为零。于是 $(\mu,L)$ 描述了当前 chunk
全部有效候选的稳定 softmax 状态。

\textbf{(b) 当前 chunk 与此前 chunk 的 softmax 状态合并。}

设处理当前 chunk 之前，已经处理的候选集合为 $\mathcal S$，
累计状态为 $(m,\ell,A)$，其中
\[
  \ell
  =
  \sum_{j\in\mathcal S}2^{z_j-m}.
\]
加入当前 chunk 后，新的最大值为
\[
  m'=\max(m,\mu).
\]
定义两个重缩放因子
\[
  \alpha=2^{m-m'},\qquad
  \gamma=2^{\mu-m'}.
\]
其中 $\alpha$ 用于重缩放此前 chunk 的结果，$\gamma$ 用于重缩放
当前 chunk 的结果。

合并后的指数和为
\[
  \boxed{
    \ell'=\alpha\ell+\gamma L
  }.
\]
此前候选的贡献满足
\[
  \alpha\ell
  =
  2^{m-m'}
  \sum_{j\in\mathcal S}2^{z_j-m}
  =
  \sum_{j\in\mathcal S}2^{z_j-m'}.
\]
当前 chunk 的贡献满足
\[
  \gamma L
  =
  2^{\mu-m'}
  \sum_{j=0}^{63}2^{z_j-\mu}
  =
  \sum_{j=0}^{63}2^{z_j-m'}.
\]
因此
\[
  \ell'
  =
  \sum_{j\in\mathcal S\cup\mathcal J_c}
  2^{z_j-m'},
\]
其中 $\mathcal J_c$ 表示当前 chunk 的有效候选集合。
这与一次性对所有已处理候选计算 softmax 分母相同。

\textbf{(c) PV 分子的分块更新。}

当前 chunk 在局部最大值 $\mu$ 基准下的 PV 分子为
\[
  B
  =
  \sum_{j\in\mathcal J_c}
  2^{z_j-\mu}V_{c,j,:}.
\]
此前 chunk 的累计 PV 分子为
\[
  A
  =
  \sum_{j\in\mathcal S}
  2^{z_j-m}V_{j,:}.
\]
为了把两部分相加，需要将它们统一到新的最大值 $m'$：
\[
  \boxed{
    A'=\alpha A+\gamma B
  }.
\]
此前候选的部分为
\[
  \alpha A
  =
  \sum_{j\in\mathcal S}
  2^{z_j-m'}V_{j,:},
\]
当前 chunk 的部分为
\[
  \gamma B
  =
  \sum_{j\in\mathcal J_c}
  2^{z_j-m'}V_{c,j,:}.
\]
所以
\[
  A'
  =
  \sum_{j\in\mathcal S\cup\mathcal J_c}
  2^{z_j-m'}V_{j,:}.
\]
合并当前 chunk 后的归一化输出为
\[
  \boxed{
    O=\frac{A'}{\ell'}
  }.
\]
展开后
\[
  O
  =
  \frac{
    \sum_{j\in\mathcal S\cup\mathcal J_c}
    2^{z_j}V_{j,:}
  }{
    \sum_{j\in\mathcal S\cup\mathcal J_c}
    2^{z_j}
  },
\]
与一次性计算全部候选的 softmax--PV 完全相同。

\textbf{(d) 从单个 head 恢复到 64 个 head 的矩阵形式。}

以上推导针对固定的一个 head。恢复全部 64 个 head 后，将各 head 的
重缩放系数组成向量
\[
  \boldsymbol{\alpha}
  =
  [\alpha_0,\ldots,\alpha_{63}]^\top
  \in\mathbb R^{64}.
\]
当前 chunk 的权重、V 和历史 PV 累加量分别为
\[
  W\in\mathbb R^{64\times64},\qquad
  V_c\in\mathbb R^{64\times512},\qquad
  A\in\mathbb R^{64\times512}.
\]
于是 64 个 head 的更新可以统一写成
\[
  \boxed{
    A'
    =
    \diag(\boldsymbol{\alpha})A+WV_c
  }.
\]
其中
\[
  \underbrace{\diag(\boldsymbol{\alpha})A}_{64\times512}
\]
表示分别重缩放 64 个 head 的历史 PV 累加量，而
\[
  \underbrace{W}_{64\times64}
  \underbrace{V_c}_{64\times512}
  \longrightarrow
  \underbrace{WV_c}_{64\times512}
\]
表示当前 chunk 对全部 64 个 head 的 PV 贡献。

non-RoPE、RoPE 两条路径及其 NVFP4/BF16 实现已在第 3 节给出，
此处不再重复。

\textbf{(e) 从分块矩阵乘法理解 PV。}

将全部候选沿候选轴分成 $C$ 个 chunk。设所有 chunk 的指数权重已经
统一到最终最大值基准，记为
\[
  G=[G_0\mid G_1\mid\cdots\mid G_{C-1}],
  \qquad
  G_c\in\mathbb R^{64\times64}.
\]
将 V 沿候选轴纵向分块：
\[
  V=
  \begin{bmatrix}
    V_0\\
    V_1\\
    \vdots\\
    V_{C-1}
  \end{bmatrix},
  \qquad
  V_c\in\mathbb R^{64\times512}.
\]
根据分块矩阵乘法，
\begin{align*}
  GV
  &=
  [G_0\mid G_1\mid\cdots\mid G_{C-1}]
  \begin{bmatrix}
    V_0\\
    V_1\\
    \vdots\\
    V_{C-1}
  \end{bmatrix}\\
  &=
  \boxed{
    \sum_{c=0}^{C-1}G_cV_c
  }.
\end{align*}
softmax 分母也可以沿相同的候选轴分块：
\[
  G\mathbf 1_{64C}
  =
  \sum_{c=0}^{C-1}G_c\mathbf 1_{64}.
\]
所以完整输出可以写成
\[
  \boxed{
  O
  =
  \diag\!\left(
    \frac{1}{
      \sum_{c=0}^{C-1}G_c\mathbf 1_{64}
    }
  \right)
  \left(
    \sum_{c=0}^{C-1}G_cV_c
  \right)
  }.
\]
这里向量的倒数按元素计算。

online softmax 不需要保存全部 $G_c$。每处理一个新 chunk，只需：

\[
  m'=\max(m,\mu),
\]
\[
  \ell'=\alpha\ell+\gamma L,
\]
\[
  A'=\alpha A+WV_c.
\]
处理完全部 chunk 后，再执行
\[
  \boxed{
    O=\frac{A}{\ell}
  },
\]
其中除法沿 head 维逐行广播，即得到与未分块 softmax--PV 相同的结果。

\section*{6\quad 一个 CTA 的执行流程}

以一个负责 64 个 head、连续若干个 chunk 的 CTA 为例，每个 chunk 有 64 个候选 KV。

\begin{enumerate}
  \item \textbf{开始时准备 Q。}
  读取这个请求的 Q：non-RoPE 部分量化成 NVFP4，RoPE 部分保留 BF16。
  初始化每个 head 的累计最大分数、指数和，以及输出累加量。

  \item \textbf{读取本轮 KV。}
  根据稀疏索引，从缓存取出本轮 64 个 KV 条目，无效候选置零并屏蔽。

  \item \textbf{计算本轮 QK 分数。}
  用 Q 与这 64 个 K 做点积，合并 non-RoPE 和 RoPE 部分的贡献，
  得到每个 head 对各候选的分数。

  \item \textbf{更新 online softmax。}
  将本轮分数与此前的 softmax 状态合并，更新累计最大值和指数和。
  若最大值变化，就相应缩放之前的输出累加量。
  同时生成本轮的指数权重，此时还不除以分母。

  \item \textbf{计算本轮 PV，累加输出。}
  non-RoPE 路径将权重量化，并将 V 按候选方向重新量化，再计算 PV；
  RoPE 路径使用 BF16。把本轮贡献加到输出累加量中。
\end{enumerate}

\textbf{每轮结束，需要保留的是：每个 head 的累计最大值、累计指数和，
以及 512 维的未归一化输出累加量，均为 FP32。}
Q 继续复用；本轮 KV、分数和权重都是临时数据，缓冲可供后续轮次使用。

重复第 2--5 步，直到该 CTA 的 chunk 全部处理完。
最后用累计分母归一化：单 CTA 直接结合 sink 写最终输出；
多个 CTA 分片时，写 BF16 中间输出和 FP32 LSE，交给后续 kernel 合并。

\end{document}